% ============================================================
% 基礎演習Ⅱ　第1回〜第5回 総まとめ（付録：使ってみたいAIプロンプト集）
%   前半「探究の基礎」5回の講義内容の総まとめと、プロンプト集・用語集。
%   授業サイトの各回ページ（_sessions/01〜05）、ワークシートと記入例、
%   AI利用ガイド、出席確認クイズ（_attendance/01〜05）の内容をもとにしている。
%   ビルド：make summary-01-05.pdf（Makefile を参照）
% ============================================================
% ============================================================
% 基礎演習Ⅱ　総まとめ冊子　共通プリアンブル
%   各冊子（summary-01-05.tex など）の冒頭で % ============================================================
% 基礎演習Ⅱ　総まとめ冊子　共通プリアンブル
%   各冊子（summary-01-05.tex など）の冒頭で % ============================================================
% 基礎演習Ⅱ　総まとめ冊子　共通プリアンブル
%   各冊子（summary-01-05.tex など）の冒頭で \input{preamble} して使う。
%   LuaLaTeX（jlreq）で組む。色と部品はワークシート（../worksheets/preamble.tex）に合わせてある。
% ============================================================
\documentclass[paper=a4paper,fontsize=10pt,line_length=48zw,number_of_lines=44,
  baselineskip=1.6zw,head_space=20mm]{jlreq}

\usepackage{xcolor}
\usepackage[most]{tcolorbox}
\usepackage{tabularray}
\usepackage{enumitem}
\usepackage{needspace}
\usepackage{fancyhdr}
\usepackage[hidelinks]{hyperref}

% ------------------------------------------------------------
% 色（サイトの assets/css/style.css、ワークシートに合わせる）
% ------------------------------------------------------------
\definecolor{ink}{HTML}{1D2330}
\definecolor{inktwo}{HTML}{4A5261}
\definecolor{inkthree}{HTML}{757D8A}
\definecolor{rulegray}{HTML}{B9BEB4}
\definecolor{papertwo}{HTML}{ECEEE8}
\definecolor{pone}{HTML}{2747C9}  \definecolor{ponetint}{HTML}{DCE2F9}
\definecolor{ptwo}{HTML}{D68A17}  \definecolor{ptwotint}{HTML}{F6E4C3}
\definecolor{bad}{HTML}{B3261E}   \definecolor{badtint}{HTML}{F9DEDA}

% ------------------------------------------------------------
% 冊子ごとに設定する情報（各冊子の冒頭で \renewcommand する）
% ------------------------------------------------------------
\newcommand*\doctitle{総まとめ}
\newcommand*\docauthor{河合 勝彦}
\newcommand*\docaffil{名古屋市立大学 経済学部}
\newcommand*\docdate{}
\AtBeginDocument{\hypersetup{pdftitle={基礎演習Ⅱ　\doctitle},pdfauthor={\docauthor}}}

% ------------------------------------------------------------
% ページスタイル
% ------------------------------------------------------------
\pagestyle{fancy}
\fancyhf{}
\renewcommand{\headrulewidth}{0pt}
\fancyfoot[L]{\sffamily\scriptsize\color{inkthree}基礎演習Ⅱ　2026年度後期　\doctitle}
\fancyfoot[R]{\sffamily\scriptsize\color{inkthree}\thepage}
\fancypagestyle{cover}{\fancyhf{}}

\setlist{nosep,leftmargin=1.5em}
\setcounter{tocdepth}{2}
\widowpenalty=10000 \clubpenalty=10000

% ------------------------------------------------------------
% 見出し
% ------------------------------------------------------------
\ModifyHeading{section}{font=\large\sffamily\bfseries\color{pone},
  before_space=1.2\baselineskip,after_space=0.5\baselineskip,lines=3,
  label_format={\thesection},after_label_space=1\zw}
\ModifyHeading{subsection}{font=\normalsize\sffamily\bfseries\color{ink},
  before_space=0.9\baselineskip,after_space=0.3\baselineskip,lines=2,
  label_format={\thesubsection},after_label_space=0.8\zw}

% 部の帯　\partband{ラベル}{題}
\newcommand\partband[2]{%
  \clearpage
  \begin{tcolorbox}[enhanced,colback=pone,colframe=pone,arc=2mm,boxrule=0pt,
      left=4mm,right=4mm,top=2.5mm,bottom=2.5mm,fontupper=\color{white}\sffamily]
    {\small #1}\par\vspace{0.6mm}{\LARGE\bfseries #2}
  \end{tcolorbox}\vspace{1mm}}

% 付録に切り替える（節番号を 付録A、A.1 にする）
\newcommand\appendixmode{%
  \setcounter{section}{0}%
  \renewcommand\thesection{付録\Alph{section}}%
  \renewcommand\thesubsection{\Alph{section}.\arabic{subsection}}%
  \renewcommand\theHsection{appendix.\Alph{section}}%
  \renewcommand\theHsubsection{appendix.\Alph{section}.\arabic{subsection}}}

% 各回の見出し　\sessionhead{回}{日付}{題}{タグ}
\newcommand\sessionhead[4]{%
  \par\Needspace{34mm}\vspace{0.9\baselineskip}%
  \phantomsection\addcontentsline{toc}{subsection}{第#1回　#3}%
  \begin{tcolorbox}[enhanced,colback=ponetint,colframe=ponetint,arc=1.5mm,boxrule=0pt,
      borderline west={1.5mm}{0pt}{pone},left=4mm,right=3mm,top=1.5mm,bottom=1.5mm,
      before skip=0pt,after skip=2.5mm]
    {\sffamily\small\color{pone}第#1回　#2\if\relax\detokenize{#4}\relax\else　\textbar　#4\fi}\par
    \vspace{0.3mm}{\sffamily\large\bfseries\color{ink}#3}
  \end{tcolorbox}}

% ------------------------------------------------------------
% 囲み
% ------------------------------------------------------------
% 要点　\begin{keybox}{見出し} ... \end{keybox}
\newtcolorbox{keybox}[1]{enhanced,colback=papertwo,colframe=papertwo,boxrule=0pt,arc=1.5mm,
  left=3.5mm,right=3mm,top=1.5mm,bottom=1.5mm,fontupper=\small,
  before upper={{\sffamily\bfseries\color{pone}#1}\par\vspace{0.4mm}},
  before skip=2mm,after skip=2.5mm}

% 注意　\begin{caution}{見出し} ... \end{caution}
\newtcolorbox{caution}[1]{enhanced,colback=badtint,colframe=badtint,boxrule=0pt,arc=1.5mm,
  left=3.5mm,right=3mm,top=1.5mm,bottom=1.5mm,fontupper=\small,
  before upper={{\sffamily\bfseries\color{bad}#1}\par\vspace{0.4mm}},
  before skip=2mm,after skip=2.5mm}

% プロンプト　\begin{promptbox}[題] ... \end{promptbox}
%   （　）の部分は自分の内容に置き換える。行は空行で区切る。
\newtcolorbox{promptbox}[1][]{enhanced,colback=papertwo,colframe=papertwo,boxrule=0pt,
  arc=1mm,borderline west={1.2mm}{0pt}{pone},left=3.5mm,right=3mm,top=1.5mm,bottom=1.8mm,
  fontupper=\sffamily\small\color{ink}\parindent0pt,before skip=1.5mm,after skip=2mm,
  title={\sffamily\scriptsize\bfseries\color{pone}プロンプト\if\relax\detokenize{#1}\relax\else　#1\fi},
  attach title to upper={\par\vspace{0.6mm}}}

% チェック項目
\newenvironment{checks}{\begin{itemize}[label={□},leftmargin=1.5em,itemsep=0.8mm]\small}{\end{itemize}}

% ------------------------------------------------------------
% プロンプト集の部品
% ------------------------------------------------------------
% 項目の見出し　\promptitem{番号}{題}{使う場面}
\newcommand\promptitem[3]{%
  \par\Needspace{48mm}\vspace{1.2\baselineskip}%
  \phantomsection\label{p:#1}%
  \noindent{\sffamily\bfseries\colorbox{pone}{\strut\color{white}\,#1\,}}\enspace
  {\sffamily\bfseries\large #2}\par
  \vspace{0.4mm}\noindent{\small\color{inktwo}使う場面：#3}\par\vspace{0.6mm}}

% 注記　\pnote{ラベル}{本文}（ラベルは「解説」「続けて聞く」「確かめる」など）
\newcommand\pnote[2]{%
  \par\vspace{0.6mm}{\small\noindent\hangindent=5.5\zw\hangafter=1
  {\sffamily\bfseries\color{pone}\makebox[5.5\zw][l]{#1}}#2\par}}

% 授業で使ったプロンプトの印
\newcommand\usedmark{{\sffamily\scriptsize\color{white}\colorbox{ptwo}{\strut 授業で使用}}}

% ------------------------------------------------------------
% 表の共通スタイル
% ------------------------------------------------------------
\NewTblrEnviron{stab}
\SetTblrInner[stab]{hlines={0.4pt,rulegray},vlines={0.4pt,rulegray},
  cells={font=\small},rowsep=2.5pt,colsep=4pt,
  row{1}={bg=ponetint,font=\sffamily\footnotesize\bfseries}}
\newcommand\hc[1]{\SetCell{bg=ponetint,font=\sffamily\footnotesize\bfseries}#1}

% 付録の小見出し（次の項目と一緒に次の頁へ送る）
\newcommand\psubsection[1]{\par\Needspace{64mm}\subsection{#1}}
 して使う。
%   LuaLaTeX（jlreq）で組む。色と部品はワークシート（../worksheets/preamble.tex）に合わせてある。
% ============================================================
\documentclass[paper=a4paper,fontsize=10pt,line_length=48zw,number_of_lines=44,
  baselineskip=1.6zw,head_space=20mm]{jlreq}

\usepackage{xcolor}
\usepackage[most]{tcolorbox}
\usepackage{tabularray}
\usepackage{enumitem}
\usepackage{needspace}
\usepackage{fancyhdr}
\usepackage[hidelinks]{hyperref}

% ------------------------------------------------------------
% 色（サイトの assets/css/style.css、ワークシートに合わせる）
% ------------------------------------------------------------
\definecolor{ink}{HTML}{1D2330}
\definecolor{inktwo}{HTML}{4A5261}
\definecolor{inkthree}{HTML}{757D8A}
\definecolor{rulegray}{HTML}{B9BEB4}
\definecolor{papertwo}{HTML}{ECEEE8}
\definecolor{pone}{HTML}{2747C9}  \definecolor{ponetint}{HTML}{DCE2F9}
\definecolor{ptwo}{HTML}{D68A17}  \definecolor{ptwotint}{HTML}{F6E4C3}
\definecolor{bad}{HTML}{B3261E}   \definecolor{badtint}{HTML}{F9DEDA}

% ------------------------------------------------------------
% 冊子ごとに設定する情報（各冊子の冒頭で \renewcommand する）
% ------------------------------------------------------------
\newcommand*\doctitle{総まとめ}
\newcommand*\docauthor{河合 勝彦}
\newcommand*\docaffil{名古屋市立大学 経済学部}
\newcommand*\docdate{}
\AtBeginDocument{\hypersetup{pdftitle={基礎演習Ⅱ　\doctitle},pdfauthor={\docauthor}}}

% ------------------------------------------------------------
% ページスタイル
% ------------------------------------------------------------
\pagestyle{fancy}
\fancyhf{}
\renewcommand{\headrulewidth}{0pt}
\fancyfoot[L]{\sffamily\scriptsize\color{inkthree}基礎演習Ⅱ　2026年度後期　\doctitle}
\fancyfoot[R]{\sffamily\scriptsize\color{inkthree}\thepage}
\fancypagestyle{cover}{\fancyhf{}}

\setlist{nosep,leftmargin=1.5em}
\setcounter{tocdepth}{2}
\widowpenalty=10000 \clubpenalty=10000

% ------------------------------------------------------------
% 見出し
% ------------------------------------------------------------
\ModifyHeading{section}{font=\large\sffamily\bfseries\color{pone},
  before_space=1.2\baselineskip,after_space=0.5\baselineskip,lines=3,
  label_format={\thesection},after_label_space=1\zw}
\ModifyHeading{subsection}{font=\normalsize\sffamily\bfseries\color{ink},
  before_space=0.9\baselineskip,after_space=0.3\baselineskip,lines=2,
  label_format={\thesubsection},after_label_space=0.8\zw}

% 部の帯　\partband{ラベル}{題}
\newcommand\partband[2]{%
  \clearpage
  \begin{tcolorbox}[enhanced,colback=pone,colframe=pone,arc=2mm,boxrule=0pt,
      left=4mm,right=4mm,top=2.5mm,bottom=2.5mm,fontupper=\color{white}\sffamily]
    {\small #1}\par\vspace{0.6mm}{\LARGE\bfseries #2}
  \end{tcolorbox}\vspace{1mm}}

% 付録に切り替える（節番号を 付録A、A.1 にする）
\newcommand\appendixmode{%
  \setcounter{section}{0}%
  \renewcommand\thesection{付録\Alph{section}}%
  \renewcommand\thesubsection{\Alph{section}.\arabic{subsection}}%
  \renewcommand\theHsection{appendix.\Alph{section}}%
  \renewcommand\theHsubsection{appendix.\Alph{section}.\arabic{subsection}}}

% 各回の見出し　\sessionhead{回}{日付}{題}{タグ}
\newcommand\sessionhead[4]{%
  \par\Needspace{34mm}\vspace{0.9\baselineskip}%
  \phantomsection\addcontentsline{toc}{subsection}{第#1回　#3}%
  \begin{tcolorbox}[enhanced,colback=ponetint,colframe=ponetint,arc=1.5mm,boxrule=0pt,
      borderline west={1.5mm}{0pt}{pone},left=4mm,right=3mm,top=1.5mm,bottom=1.5mm,
      before skip=0pt,after skip=2.5mm]
    {\sffamily\small\color{pone}第#1回　#2\if\relax\detokenize{#4}\relax\else　\textbar　#4\fi}\par
    \vspace{0.3mm}{\sffamily\large\bfseries\color{ink}#3}
  \end{tcolorbox}}

% ------------------------------------------------------------
% 囲み
% ------------------------------------------------------------
% 要点　\begin{keybox}{見出し} ... \end{keybox}
\newtcolorbox{keybox}[1]{enhanced,colback=papertwo,colframe=papertwo,boxrule=0pt,arc=1.5mm,
  left=3.5mm,right=3mm,top=1.5mm,bottom=1.5mm,fontupper=\small,
  before upper={{\sffamily\bfseries\color{pone}#1}\par\vspace{0.4mm}},
  before skip=2mm,after skip=2.5mm}

% 注意　\begin{caution}{見出し} ... \end{caution}
\newtcolorbox{caution}[1]{enhanced,colback=badtint,colframe=badtint,boxrule=0pt,arc=1.5mm,
  left=3.5mm,right=3mm,top=1.5mm,bottom=1.5mm,fontupper=\small,
  before upper={{\sffamily\bfseries\color{bad}#1}\par\vspace{0.4mm}},
  before skip=2mm,after skip=2.5mm}

% プロンプト　\begin{promptbox}[題] ... \end{promptbox}
%   （　）の部分は自分の内容に置き換える。行は空行で区切る。
\newtcolorbox{promptbox}[1][]{enhanced,colback=papertwo,colframe=papertwo,boxrule=0pt,
  arc=1mm,borderline west={1.2mm}{0pt}{pone},left=3.5mm,right=3mm,top=1.5mm,bottom=1.8mm,
  fontupper=\sffamily\small\color{ink}\parindent0pt,before skip=1.5mm,after skip=2mm,
  title={\sffamily\scriptsize\bfseries\color{pone}プロンプト\if\relax\detokenize{#1}\relax\else　#1\fi},
  attach title to upper={\par\vspace{0.6mm}}}

% チェック項目
\newenvironment{checks}{\begin{itemize}[label={□},leftmargin=1.5em,itemsep=0.8mm]\small}{\end{itemize}}

% ------------------------------------------------------------
% プロンプト集の部品
% ------------------------------------------------------------
% 項目の見出し　\promptitem{番号}{題}{使う場面}
\newcommand\promptitem[3]{%
  \par\Needspace{48mm}\vspace{1.2\baselineskip}%
  \phantomsection\label{p:#1}%
  \noindent{\sffamily\bfseries\colorbox{pone}{\strut\color{white}\,#1\,}}\enspace
  {\sffamily\bfseries\large #2}\par
  \vspace{0.4mm}\noindent{\small\color{inktwo}使う場面：#3}\par\vspace{0.6mm}}

% 注記　\pnote{ラベル}{本文}（ラベルは「解説」「続けて聞く」「確かめる」など）
\newcommand\pnote[2]{%
  \par\vspace{0.6mm}{\small\noindent\hangindent=5.5\zw\hangafter=1
  {\sffamily\bfseries\color{pone}\makebox[5.5\zw][l]{#1}}#2\par}}

% 授業で使ったプロンプトの印
\newcommand\usedmark{{\sffamily\scriptsize\color{white}\colorbox{ptwo}{\strut 授業で使用}}}

% ------------------------------------------------------------
% 表の共通スタイル
% ------------------------------------------------------------
\NewTblrEnviron{stab}
\SetTblrInner[stab]{hlines={0.4pt,rulegray},vlines={0.4pt,rulegray},
  cells={font=\small},rowsep=2.5pt,colsep=4pt,
  row{1}={bg=ponetint,font=\sffamily\footnotesize\bfseries}}
\newcommand\hc[1]{\SetCell{bg=ponetint,font=\sffamily\footnotesize\bfseries}#1}

% 付録の小見出し（次の項目と一緒に次の頁へ送る）
\newcommand\psubsection[1]{\par\Needspace{64mm}\subsection{#1}}
 して使う。
%   LuaLaTeX（jlreq）で組む。色と部品はワークシート（../worksheets/preamble.tex）に合わせてある。
% ============================================================
\documentclass[paper=a4paper,fontsize=10pt,line_length=48zw,number_of_lines=44,
  baselineskip=1.6zw,head_space=20mm]{jlreq}

\usepackage{xcolor}
\usepackage[most]{tcolorbox}
\usepackage{tabularray}
\usepackage{enumitem}
\usepackage{needspace}
\usepackage{fancyhdr}
\usepackage[hidelinks]{hyperref}

% ------------------------------------------------------------
% 色（サイトの assets/css/style.css、ワークシートに合わせる）
% ------------------------------------------------------------
\definecolor{ink}{HTML}{1D2330}
\definecolor{inktwo}{HTML}{4A5261}
\definecolor{inkthree}{HTML}{757D8A}
\definecolor{rulegray}{HTML}{B9BEB4}
\definecolor{papertwo}{HTML}{ECEEE8}
\definecolor{pone}{HTML}{2747C9}  \definecolor{ponetint}{HTML}{DCE2F9}
\definecolor{ptwo}{HTML}{D68A17}  \definecolor{ptwotint}{HTML}{F6E4C3}
\definecolor{bad}{HTML}{B3261E}   \definecolor{badtint}{HTML}{F9DEDA}

% ------------------------------------------------------------
% 冊子ごとに設定する情報（各冊子の冒頭で \renewcommand する）
% ------------------------------------------------------------
\newcommand*\doctitle{総まとめ}
\newcommand*\docauthor{河合 勝彦}
\newcommand*\docaffil{名古屋市立大学 経済学部}
\newcommand*\docdate{}
\AtBeginDocument{\hypersetup{pdftitle={基礎演習Ⅱ　\doctitle},pdfauthor={\docauthor}}}

% ------------------------------------------------------------
% ページスタイル
% ------------------------------------------------------------
\pagestyle{fancy}
\fancyhf{}
\renewcommand{\headrulewidth}{0pt}
\fancyfoot[L]{\sffamily\scriptsize\color{inkthree}基礎演習Ⅱ　2026年度後期　\doctitle}
\fancyfoot[R]{\sffamily\scriptsize\color{inkthree}\thepage}
\fancypagestyle{cover}{\fancyhf{}}

\setlist{nosep,leftmargin=1.5em}
\setcounter{tocdepth}{2}
\widowpenalty=10000 \clubpenalty=10000

% ------------------------------------------------------------
% 見出し
% ------------------------------------------------------------
\ModifyHeading{section}{font=\large\sffamily\bfseries\color{pone},
  before_space=1.2\baselineskip,after_space=0.5\baselineskip,lines=3,
  label_format={\thesection},after_label_space=1\zw}
\ModifyHeading{subsection}{font=\normalsize\sffamily\bfseries\color{ink},
  before_space=0.9\baselineskip,after_space=0.3\baselineskip,lines=2,
  label_format={\thesubsection},after_label_space=0.8\zw}

% 部の帯　\partband{ラベル}{題}
\newcommand\partband[2]{%
  \clearpage
  \begin{tcolorbox}[enhanced,colback=pone,colframe=pone,arc=2mm,boxrule=0pt,
      left=4mm,right=4mm,top=2.5mm,bottom=2.5mm,fontupper=\color{white}\sffamily]
    {\small #1}\par\vspace{0.6mm}{\LARGE\bfseries #2}
  \end{tcolorbox}\vspace{1mm}}

% 付録に切り替える（節番号を 付録A、A.1 にする）
\newcommand\appendixmode{%
  \setcounter{section}{0}%
  \renewcommand\thesection{付録\Alph{section}}%
  \renewcommand\thesubsection{\Alph{section}.\arabic{subsection}}%
  \renewcommand\theHsection{appendix.\Alph{section}}%
  \renewcommand\theHsubsection{appendix.\Alph{section}.\arabic{subsection}}}

% 各回の見出し　\sessionhead{回}{日付}{題}{タグ}
\newcommand\sessionhead[4]{%
  \par\Needspace{34mm}\vspace{0.9\baselineskip}%
  \phantomsection\addcontentsline{toc}{subsection}{第#1回　#3}%
  \begin{tcolorbox}[enhanced,colback=ponetint,colframe=ponetint,arc=1.5mm,boxrule=0pt,
      borderline west={1.5mm}{0pt}{pone},left=4mm,right=3mm,top=1.5mm,bottom=1.5mm,
      before skip=0pt,after skip=2.5mm]
    {\sffamily\small\color{pone}第#1回　#2\if\relax\detokenize{#4}\relax\else　\textbar　#4\fi}\par
    \vspace{0.3mm}{\sffamily\large\bfseries\color{ink}#3}
  \end{tcolorbox}}

% ------------------------------------------------------------
% 囲み
% ------------------------------------------------------------
% 要点　\begin{keybox}{見出し} ... \end{keybox}
\newtcolorbox{keybox}[1]{enhanced,colback=papertwo,colframe=papertwo,boxrule=0pt,arc=1.5mm,
  left=3.5mm,right=3mm,top=1.5mm,bottom=1.5mm,fontupper=\small,
  before upper={{\sffamily\bfseries\color{pone}#1}\par\vspace{0.4mm}},
  before skip=2mm,after skip=2.5mm}

% 注意　\begin{caution}{見出し} ... \end{caution}
\newtcolorbox{caution}[1]{enhanced,colback=badtint,colframe=badtint,boxrule=0pt,arc=1.5mm,
  left=3.5mm,right=3mm,top=1.5mm,bottom=1.5mm,fontupper=\small,
  before upper={{\sffamily\bfseries\color{bad}#1}\par\vspace{0.4mm}},
  before skip=2mm,after skip=2.5mm}

% プロンプト　\begin{promptbox}[題] ... \end{promptbox}
%   （　）の部分は自分の内容に置き換える。行は空行で区切る。
\newtcolorbox{promptbox}[1][]{enhanced,colback=papertwo,colframe=papertwo,boxrule=0pt,
  arc=1mm,borderline west={1.2mm}{0pt}{pone},left=3.5mm,right=3mm,top=1.5mm,bottom=1.8mm,
  fontupper=\sffamily\small\color{ink}\parindent0pt,before skip=1.5mm,after skip=2mm,
  title={\sffamily\scriptsize\bfseries\color{pone}プロンプト\if\relax\detokenize{#1}\relax\else　#1\fi},
  attach title to upper={\par\vspace{0.6mm}}}

% チェック項目
\newenvironment{checks}{\begin{itemize}[label={□},leftmargin=1.5em,itemsep=0.8mm]\small}{\end{itemize}}

% ------------------------------------------------------------
% プロンプト集の部品
% ------------------------------------------------------------
% 項目の見出し　\promptitem{番号}{題}{使う場面}
\newcommand\promptitem[3]{%
  \par\Needspace{48mm}\vspace{1.2\baselineskip}%
  \phantomsection\label{p:#1}%
  \noindent{\sffamily\bfseries\colorbox{pone}{\strut\color{white}\,#1\,}}\enspace
  {\sffamily\bfseries\large #2}\par
  \vspace{0.4mm}\noindent{\small\color{inktwo}使う場面：#3}\par\vspace{0.6mm}}

% 注記　\pnote{ラベル}{本文}（ラベルは「解説」「続けて聞く」「確かめる」など）
\newcommand\pnote[2]{%
  \par\vspace{0.6mm}{\small\noindent\hangindent=5.5\zw\hangafter=1
  {\sffamily\bfseries\color{pone}\makebox[5.5\zw][l]{#1}}#2\par}}

% 授業で使ったプロンプトの印
\newcommand\usedmark{{\sffamily\scriptsize\color{white}\colorbox{ptwo}{\strut 授業で使用}}}

% ------------------------------------------------------------
% 表の共通スタイル
% ------------------------------------------------------------
\NewTblrEnviron{stab}
\SetTblrInner[stab]{hlines={0.4pt,rulegray},vlines={0.4pt,rulegray},
  cells={font=\small},rowsep=2.5pt,colsep=4pt,
  row{1}={bg=ponetint,font=\sffamily\footnotesize\bfseries}}
\newcommand\hc[1]{\SetCell{bg=ponetint,font=\sffamily\footnotesize\bfseries}#1}

% 付録の小見出し（次の項目と一緒に次の頁へ送る）
\newcommand\psubsection[1]{\par\Needspace{64mm}\subsection{#1}}

\renewcommand*\doctitle{第1回〜第5回 総まとめ}
\renewcommand*\docdate{2026年10月28日}

\begin{document}

% ============================================================
% 表紙
% ============================================================
\thispagestyle{cover}
\vspace*{8mm}
{\sffamily\color{pone}\large 名古屋市立大学 経済学部　2026年度後期　水曜3限}\par
\vspace{3mm}
{\sffamily\bfseries\fontsize{30pt}{38pt}\selectfont 基礎演習Ⅱ}\par
\vspace{2mm}
{\sffamily\fontsize{20pt}{26pt}\selectfont 第1回〜第5回　総まとめ}\par
\vspace{2.5mm}
{\sffamily\color{inktwo}生成AIを「思考のパートナー」にする探究の基礎　／　付録：使ってみたいAIプロンプト集}\par
\vspace{4mm}
\noindent{\color{inktwo}\rule{\linewidth}{0.6pt}}\par
\vspace{3mm}
{\sffamily\small \docauthor（\docaffil）\hfill\docdate}\par
\vspace{5mm}

\begin{tcolorbox}[enhanced,colback=papertwo,colframe=papertwo,boxrule=0pt,arc=1.5mm,
    left=4mm,right=4mm,top=3mm,bottom=3mm,fontupper=\small]
{\sffamily\bfseries この冊子の使い方}\par\vspace{1mm}
\begin{itemize}[label=・,itemsep=0.6mm]
  \item 第1部は、前半5回（探究の基礎）で学んだことの総まとめです。各回の要点と、5回を貫く考え方を3ページ程度にまとめました。授業サイトの各回のページ、ワークシート、AI利用ガイド、出席確認クイズの内容をもとにしています。
  \item 付録Aは「使ってみたいAIプロンプト集」です。授業で使ったプロンプトに、新しいプロンプトを加えて28個を集め、それぞれに解説、続けて聞くときの例、使う前に確かめることを添えました。（　）の部分は自分の内容に置き換えて使ってください。
  \item 付録Bは、出席確認クイズで扱ったAIの用語の一覧です。
  \item 使うAIは大学が提供する Copilot です。出力は素材であって完成品ではありません。事実、数値、文献は一次資料で確かめ、使った場面はAI活用記録に書きます。
\end{itemize}
\end{tcolorbox}

\vspace{4mm}
{\sffamily\bfseries 目次}\par\vspace{0.5mm}
{\footnotesize\makeatletter\@starttoc{toc}\makeatother}

% ============================================================
% 第1部　総まとめ
% ============================================================
\partband{第1部}{前半5回の総まとめ}

\section{前半5回で何をしたか}

前半の5回は「探究の基礎」のフェーズでした。ねらいは、生成AIを「思考のパートナー」として使いながら、自分の関心を、学術的に答えられる問い（リサーチ・クエスチョン）に育てることです。流れは一本道でした。第1回で持ち帰った「関心のあるテーマ」が、第2回で仮の問いになり、第3回で先行研究と出会い、第4回で1本の論文を深く読むことで問いの隙間が見え、第5回でAIの批判を受けて確定しました。

5回を通じて、AIの役割は毎回変わりました。最初は「嘘をつく相手」として検証の必要性を教え、次に「ブレストの相手」として視点を広げ、「探す場所の案内役」「読む前の地図と議論の相手」を経て、最後は「厳しい査読者」になりました。一方で、問いを選ぶこと、事実を確かめること、自分の言葉で書くことは、どの回でも自分の仕事でした。

\noindent\begin{stab}{colspec={Q[c,wd=8mm,m] Q[l,wd=30mm,m] X[2,l,m] X[3,l,m] X[2,l,m]},cells={font=\footnotesize},rowsep=2pt}
  回 & 題（日付） & AIの役割 & 身につけたこと & 持ち帰ったもの \\
  1 & ガイダンス\newline（9月30日） & 嘘をつく相手（ハルシネーションの体験） & 授業の仕組みと評価、AI利用のルール、出力を使う前のチェック & 関心のあるテーマ3つと、気になる理由 \\
  2 & 問いの立て方\newline（10月7日） & ブレストの相手（6つの視点） & テーマと問いの違い、良い問いの3条件、採点して絞る & 仮のリサーチ・クエスチョンと選んだ理由 \\
  3 & 先行研究の探索\newline（10月14日） & 探す場所の案内役（分野とキーワード） & 学術データベースの検索、文献の実在と内容の確認、文献リストの6項目 & 実在を確認した文献5件、精読する1件 \\
  4 & 文献講読と批判的思考\newline（10月21日） & 読む前の地図、議論の相手 & 論文の構造、AI要約のずれの見つけ方、反論に原文の根拠で答える & 主張・根拠・限界の要約、未解決の要素 \\
  5 & リサーチ・クエスチョンの確定\newline（10月28日） & 厳しい査読者 & 4観点での検証、補強と修正の使い分け、1枚にまとめて発表 & 確定した問い、確定レポート \\
\end{stab}

\section{各回の要点}

\sessionhead{1}{9月30日（水）}{ガイダンス}{}

授業は「探究の基礎（第1〜5回）」と「実践プロジェクト（第6〜13回）」の二部構成で、第14・15回はキャリア講習です。成績は提出物50\%と発表50\%で評価します。前半のリサーチ・クエスチョン確定レポート、後半の店舗設計提案、最後の振り返りレポートのすべてに、AI活用記録を添付します。

この授業でのAI利用の約束は、AI利用ガイドの4つの考え方に集約されます。

\begin{keybox}{AI利用の4つの考え方}
\begin{enumerate}[label=\arabic*.,leftmargin=1.6em,itemsep=0.3mm]
  \item \textbf{問いは自分が持つ。}テーマを決め、何を知りたいかを言葉にするのは自分の仕事。AIは視点を広げる相手
  \item \textbf{出力は素材であって完成品ではない。}読み、検証し、自分の言葉で書き直す
  \item \textbf{事実は一次資料で確かめる。}数値、固有名詞、文献情報は、元の資料（論文、統計、公式サイト）で確認する
  \item \textbf{使ったことを隠さない。}どこでどうAIを使ったかを記録し、提出物に明記する。透明性は評価の対象
\end{enumerate}
\end{keybox}

後半のミニワーク「AIに嘘をつかせる」では、名古屋市立大学経済学部の歴史を参考文献付きで説明させ、挙がった文献を図書館の蔵書検索や CiNii Research で確かめました。実在しない文献が「実在しそうな体裁」で出てくること、自分がよく知っている話題（出身地、部活）でも年号や固有名詞がずれることを、自分の目で確かめたはずです。これがハルシネーションです。大規模言語モデル（LLM）は、文脈から「次に来る単語」を確率的に予測して文章を作っているだけで、意味を理解して調べているわけではありません。だから、もっともらしい誤りが自信のある口調で出てきます。

この体験から整理した「出力を使う前のチェック」は、その後の4回すべてで使いました。固有名詞は実在するか、数値には出典があるか、根拠のない断定になっていないか、反対の見解や例外はないか、自分で説明できるか、の5点です。

\sessionhead{2}{10月7日（水）}{問いの立て方}{AIとブレスト}

「テーマ」は話題の範囲で、「問い」は答えを求める疑問文です。「地方の自動車販売」はテーマ、「なぜ若い世代の来店は減っているのか」が問いです。「〜について」で終わる文は、まだ問いになっていません。良い問いの3条件は、①調べて答えられる、②範囲が広すぎず狭すぎない、③答える意義がある、の3つでした。

ワークでは、関心テーマをAIに伝え、経済学、経営学、社会学、心理学、都市計画、テクノロジーの6つの視点から問いの候補を出させ、3条件で採点して上位3つに絞り、ペアで「それはどうやって調べるの」と質問し合って、仮の問いを1つ決めました。例えば、「若い世代と自動車ディーラー」というテーマからは、「名古屋都市圏の20代にとって、オンラインでの情報収集は、ディーラーへの来店のしかたをどう変えたのか」という仮の問いが作れます。採点の数字そのものより、点数の理由を自分で説明できることが大切です。

AIとのブレストで覚えておきたいのは、AIは視点を広げる相手であって、選ぶのは自分だということです。採用しなかった案と、その理由をメモしておく練習が、第5回のレポートに添えるAI活用記録につながりました。出席クイズでは、プロンプト（AIに与える指示や質問の文章）の質が出力の質を決めること、この授業で意識する4要素が「役割」「目的と背景」「条件」「出力の形式」であること、同じプロンプトでも出力が毎回少し違うのは次の単語を確率にもとづいて選んでいるからであることを確かめました。

\sessionhead{3}{10月14日（水）}{先行研究の探索}{AI検索の活用}

この回の役割分担は明確です。AIには「どこを探すか」の見当（関連する学術分野と、日本語と英語の検索キーワード）をつけてもらい、文献そのものは CiNii Research、Google Scholar、J-STAGE、図書館の蔵書検索で探して、実在と内容を確かめます。AIに「文献を挙げて」と頼むワークでは、実在しない文献だけでなく、「実在するが内容が違う」「調査は実在するが、AIの示した数値が報告書にない」という、より見つけにくいずれがあることを学びました。AIに「本当ですか」と聞き直しても検証にはなりません。確認はデータベースと原文で行います。

文献リストには、著者、発行年、題名、掲載誌または出版社、巻号と頁、DOIまたはURLの6項目を揃えます。欠けやすいのは年とURLです。検索した日と件数も記録しておくと、あとで再現できます。文献は、公的統計や白書だけでなく、CiNii Research や J-STAGE で見つけた学術論文を少なくとも2件は含めます。

クイズで扱ったのは、知識のカットオフ（学習データはある時点までで、それ以降の出来事を知らない）と、検索や文書と連携して回答するRAG（検索拡張生成）です。RAGは最新の情報を扱えますが、検索結果の読み違いや誤った要約は起こります。回答に出典のURLが付いていても、必ず開いて、その内容が本当に書かれているかを確かめます。

\sessionhead{4}{10月21日（水）}{文献講読と批判的思考}{AI要約と対話}

論文は「問い、方法、結果、考察、限界」の構造で読みます。AIの要約は読む前の地図として使い、根拠は必ず原文で確かめます。架空論文Ｘ（または各自の文献）で行った照合では、AI要約に典型的な4つのずれがあることがわかりました。問いを広げる（「なぜ行かなくなったか」は論文の問いではない）、条件を落とす（「購入経験者に限定」が抜ける）、結果を言い過ぎる（「店に行かない」）、限界の節を見落とす、の4つです。原文の頁を書きながら照合すると、このずれに気づけます。

次に、AIに「最も強い反論」を3つ出させ、原文の根拠で答えられるかを○△×で判定しました。答えられない反論が出るのは悪いことではありません。それが論文の範囲の外だと言えれば十分で、「論文が答えていないこと」と「自分の問いとの隙間」が、次の回で問いを磨く材料になります。架空論文Ｘの例では、来店の「目的」の変化が調べられていないことが、問いを磨く材料になりました。

論文の全文を入力するときは、大学のAI利用規程と、その資料の利用条件を確認します。要旨や必要な部分だけで足りることが多いです。クイズでは、学習データの偏りがそのまま出力に出るバイアス、モデルが文章を扱う単位のトークン、一度に扱える長さの上限であるコンテキストウィンドウを確かめました。長い文書を貼ると前半が考慮されないことがあるのは、このためです。

\sessionhead{5}{10月28日（水）}{リサーチ・クエスチョンの確定}{AIによる批判的検証}

仮の問いと背景、主な先行研究3件をAIに渡し、「厳しい査読者」として、①半年で調べて答えられるか、②先行研究に対する新しさがあるか、③範囲は適切か、④用語は明確か、の4観点から批判させました。指摘には、答えられるものは「補強」で、答えられないものは問いの「修正」で応じます。すべての指摘を受け入れる必要はなく、根拠を示して反論できた点はむしろ評価されます。

例えば、「来店のしかたをどう変えたのか」という過去との比較を含む問いは、半年では答えにくいと指摘されます。これを「車の購入を検討した名古屋都市圏の20代は、来店前にオンラインで何を調べ、何を目的にディーラーへ来店しているのか」という、現状を記述する問いに直せば、答えられる問いになります。対象も「20代全体」から「この1年に購入を検討した20代」に絞り、あいまいな「来店のしかた」を「回数」と「目的」に分けて定義します。修正前と修正後の問いを並べて記録することが、レポートのAI活用記録の中心になります。

確定した問いは「問い、背景と意義、先行研究でわかっていること、調べ方の見通し」の4項目で1枚にまとめ、1人3分で発表しました。「20人のアンケートで20代と言えるのか」という質問には、一般化はできないので傾向を探る調査と位置づけ、限界として書く、と答えるのが筋です。クイズでは、LLMが計算や論理パズルで間違える理由（言語のパターンから「もっともらしい続き」を予測しているだけ）、口調の自信と内容の正しさは無関係であること、機械学習と教師あり学習の意味、出典を確認できない数値は使わないことを確かめました。

\section{5回を貫く3つの習慣}

\subsection{分担：広げるのはAI、選ぶ・確かめる・書くのは自分}

候補を出す、視点を増やす、反論を挙げる、難しい文章を言い換える、といった「広げる」作業はAIが得意です。何を問うか、どの案を採用するか、それは事実か、自分の言葉でどう書くか、という「選ぶ・確かめる・書く」はAIに渡しません。自分の考えを先に書いてからAIに批評させると、AIの案に引きずられにくくなります。良い出力があっても、「この案の弱点は」と聞いてから採用します。

\subsection{検証：固有名詞、数値、文献、URLは一次資料で}

\noindent\begin{stab}{colspec={Q[l,wd=28mm,m] X[l,m]},cells={font=\footnotesize},rowsep=2pt}
  確かめるもの & 確かめ方 \\
  文献 & CiNii Research、Google Scholar、J-STAGE、図書館の蔵書検索で、実在と内容の両方を確認する。確認できないものは使わない \\
  数値、統計 & 公的統計や報告書の該当ページで、同じ数値であることを確かめる。確認できない数値は使わない \\
  固有名詞、年 & 公式サイトや原資料で確認する。自分がよく知っている話題でもずれることがある \\
  出典のURL & 必ず開き、そのページに回答の内容が本当に書かれているかを見る \\
  主張 & 「根拠は」「出典は」「反対の見解は」「例外は」と尋ね、示された資料を自分で確認する。自分で説明できない内容は使わない \\
\end{stab}

\subsection{記録：AI活用記録と、本文での明記}

提出物には、使用したツール、使用した場面、主なプロンプト（1〜3個）、出力をどう扱ったか（採用、修正、不採用と理由）、検証したこと（何で確かめたか）の5項目からなるAI活用記録を添えます。本文の冒頭か末尾には、「本レポートの作成にあたり、問いの候補出しと構成案の検討に大学が提供する Copilot を使用した。出力はすべて筆者が一次資料に基づいて検証し、文章は筆者自身が執筆した」のような一文を入れます。AIの出力そのものを引用する場合は、プロンプトと日付がわかる形で示します。

\section{後半へのつながり}

第6回からは、企業の課題に取り組む「店舗設計提案プロジェクト」が始まります。ペルソナを作る、コンセプトを出す、提案の弱点を指摘させる、という作業は、前半で身につけた「広げる、絞る、批判させる、記録する」の型を、グループで使うものです。前半との違いは、企業から提供された非公開の資料やデータをAIに入力しない、SNSや外部に書かない、グループ外に共有しない、という取り扱いが加わることです。第6回で確認します。

\section{提出物（前半）のチェックリスト}

\begin{keybox}{リサーチ・クエスチョン確定レポートに含めるもの}
\begin{checks}
  \item 確定したリサーチ・クエスチョンと、選んだ理由
  \item 背景と意義
  \item 先行研究の整理（実在を確認した文献リスト付き。著者、発行年、題名、掲載誌または出版社、巻号と頁、DOIまたはURL）
  \item 調べ方の見通し（わかることと、わからないこと）
  \item AI活用記録（AI利用ガイドの形式。修正前と修正後の問いを含める）
  \item 本文の冒頭または末尾に、AI利用を明記する一文
\end{checks}
書式と締切は授業内で案内します。評価で見るのは、問いと視点の独自性、根拠の確かさと引用の正確さ、論理構成とアカデミック・ライティングの作法、AI活用の透明性と協働から得た洞察の4点です。
\end{keybox}

% ============================================================
% 付録A　プロンプト集
% ============================================================
\partband{付録}{使ってみたいAIプロンプト集と用語集}
\appendixmode

\section{使ってみたいAIプロンプト集}

授業で使ったプロンプト（「授業で使用」の印があるもの）に、新しいプロンプトを加えて、場面ごとに28個を集めました。どれも、この授業の「プロンプトの4要素」でできています。（　）の部分は自分の内容に置き換えてください。文の形はそのまま使う必要はなく、自分の状況に合わせて条件を足したり削ったりしてかまいません。

\noindent\begin{minipage}[t]{0.485\linewidth}
\noindent\begin{stab}{colspec={Q[c,wd=7mm,m] X[l,m] Q[c,wd=7mm,m]},cells={font=\footnotesize},rowsep=1.8pt,colsep=3pt}
  番号 & プロンプト & 頁 \\
  01 & 4要素の基本テンプレート & \pageref{p:01} \\
  02 & 対話を深める「続けて聞く」ことば & \pageref{p:02} \\
  03 & 自分の考えを先に書いてから批評させる & \pageref{p:03} \\
  04 & 視点を広げる & \pageref{p:04} \\
  05 & 問いを絞る & \pageref{p:05} \\
  06 & 「なぜ気になるのか」から問いを育てる & \pageref{p:06} \\
  07 & 問いの形を整える & \pageref{p:07} \\
  08 & 探す場所の見当をつける & \pageref{p:08} \\
  09 & 要旨から関連の深さを判断する & \pageref{p:09} \\
  10 & 検索式を組み立てる & \pageref{p:10} \\
  11 & 公的統計や白書のあたりをつける & \pageref{p:11} \\
  12 & 自分で確認した文献情報を整形する & \pageref{p:12} \\
  13 & 構造をつかむ & \pageref{p:13} \\
  14 & 難しい段落をかみ砕く & \pageref{p:14} \\
\end{stab}
\end{minipage}\hfill
\begin{minipage}[t]{0.485\linewidth}
\noindent\begin{stab}{colspec={Q[c,wd=7mm,m] X[l,m] Q[c,wd=7mm,m]},cells={font=\footnotesize},rowsep=1.8pt,colsep=3pt}
  番号 & プロンプト & 頁 \\
  15 & 主張と根拠の対応を確かめる & \pageref{p:15} \\
  16 & 反論を引き出す & \pageref{p:16} \\
  17 & 別の解釈を探す & \pageref{p:17} \\
  18 & 論文が答えていないことを探す & \pageref{p:18} \\
  19 & 査読者として批判させる & \pageref{p:19} \\
  20 & 反論に答える練習 & \pageref{p:20} \\
  21 & 調べ方の見通しを立てる & \pageref{p:21} \\
  22 & 発表の想定問答 & \pageref{p:22} \\
  23 & 構成案を一緒に作る & \pageref{p:23} \\
  24 & 自分の文章を推敲させる & \pageref{p:24} \\
  25 & AI活用記録の下書きを整える & \pageref{p:25} \\
  26 & 根拠と出典を尋ねる & \pageref{p:26} \\
  27 & 確かなこと、推測、わからないことに分けさせる & \pageref{p:27} \\
  28 & 反対の見解と例外を探す & \pageref{p:28} \\
\end{stab}
\end{minipage}

\begin{keybox}{プロンプトの4要素}
\begin{enumerate}[label=\arabic*.,leftmargin=1.6em,itemsep=0.3mm]
  \item \textbf{役割}　AIに立ってほしい立場。「あなたは経営学を専門とする大学教員です」「あなたは厳しい査読者です」
  \item \textbf{目的と背景}　何のために、どんな状況で必要か。「2年生の演習で、リサーチ・クエスチョンを確定するレポートを書いています」
  \item \textbf{条件}　範囲、制約、避けてほしいこと。「日本国内の事例に限る」「文献名は挙げないでください」
  \item \textbf{出力の形式}　箇条書き、表、文字数、案の数。「300字以内」「10案を表で」
\end{enumerate}
\end{keybox}

どのプロンプトにも共通する使い方が3つあります。一度で終わらせず、「なぜ」「他には」「反対の立場だと」と続けること。出力は素材として扱い、固有名詞、数値、文献は一次資料で確かめること。使った場面と、採用・不採用の判断をAI活用記録に残すことです。

% ------------------------------------------------------------
\psubsection{基本の型と、対話の続け方}

\promptitem{01}{4要素の基本テンプレート}{どの場面でも。プロンプトを一から書くとき}
\begin{promptbox}
あなたは（役割。例：経営学を専門とする大学教員）です。

私は経済学部の2年生で、（目的と背景。例：リサーチ・クエスチョンを確定するレポートを書いています）。

（依頼。例：次の問いの弱点を指摘してください。）

条件：（例：日本国内の事例に限る／文献名や著者名は挙げない／専門用語には短い説明を付ける）

出力の形式：（例：箇条書きで5つ、各2行以内／表で／300字以内）

内容：（問い、文章、データなど）
\end{promptbox}
\pnote{解説}{役割、目的と背景、条件、出力の形式を、この順に書く型です。背景が1行あるだけで出力の的が絞られます。「文献名や著者名は挙げない」のように避けてほしいことを条件に書くと、ハルシネーションの入口を減らせます。出力の形式を決めておくと、複数の案を比べやすくなります。}
\pnote{続けて聞く}{出力が抽象的なときは、具体例、数値、条件を足して聞き直します。「名古屋都市圏に限定して」「20代の購入者に絞って」のように、範囲を狭めると具体的になります。}
\pnote{確かめる}{役割を与えても、AIがその専門家になるわけではありません。内容の正しさは変わらないので、固有名詞と数値は一次資料で確かめます。}

\promptitem{02}{対話を深める「続けて聞く」ことば}{最初の出力が返ってきたあと}
\begin{promptbox}
・なぜそう言えるのですか。根拠を示してください。

・他にはありませんか。今の案と重ならないものを3つ追加してください。

・反対の立場の人なら、この案にどう反論しますか。

・この案の弱点はどこですか。最も大きいものから順に挙げてください。

・抽象的です。具体例を1つ、数値や条件を含めて示してください。

・「（用語）」はどういう意味で使っていますか。定義してください。

・今の説明を、専門知識のない人にもわかるように100字で言い直してください。
\end{promptbox}
\pnote{解説}{AIとの対話は、一度で終わらせないのがコツです。最初の出力は平均的で無難なものになりがちで、「なぜ」「他には」「反対の立場だと」と続けることで、自分の発想にない方向が出てきます。良い出力があっても、「この案の弱点は」と聞いてから採用します。}
\pnote{確かめる}{「根拠」として挙がった文献や数値も検証の対象です。根拠を尋ねたからといって、答えが正しくなるわけではありません。}

\promptitem{03}{自分の考えを先に書いてから批評させる}{自分の案があるとき。AIの案に引きずられたくないとき}
\begin{promptbox}
まず私の考えを示します。賛成や称賛はせず、弱点と見落としだけを指摘してください。

私の考え：（問い、主張、提案を自分の言葉で書く）

指摘は、重要なものから順に3つ。それぞれ、なぜ問題か、どう直せるかを1行ずつ添えてください。
\end{promptbox}
\pnote{解説}{AIは利用者の意見に同意しやすい傾向があります。「賛成しない」「称賛しない」と条件に書くと、批判が出やすくなります。自分の考えを先に文章にしておくと、AIの案に自分の考えが置き換わってしまうことを防げます。AI利用ガイドの「自分の考えを先に書いてから、AIに批評させる」の実践です。}
\pnote{続けて聞く}{「指摘1は、私の問いを変えずに直せますか」「この指摘に反論するとしたら、どんな根拠が要りますか」}
\pnote{確かめる}{指摘のうち、どれを受け入れ、どれを受け入れなかったかと、その理由をAI活用記録に書きます。}

% ------------------------------------------------------------
\psubsection{問いを立てる（第2回）}

\promptitem{04}{視点を広げる\quad\usedmark}{第2回。関心テーマから問いの候補を出すとき}
\begin{promptbox}
私は経済学部の2年生です。「（あなたのテーマ）」に関心があります。

このテーマについて、経済学、経営学、社会学、心理学、都市計画、テクノロジーの6つの視点から、それぞれ2つずつ研究の問いを提案してください。

各問いは「なぜ〜なのか」「〜は〜にどのように影響するのか」の形で、1行ずつ箇条書きにしてください。
\end{promptbox}
\pnote{解説}{視点を6つ指定することで、自分の発想にはない方向の問いが出てきます。「なぜ〜なのか」と形を決めておくと、出力が「〜について」ではなく疑問文になります。例えば「若い世代と自動車ディーラー」というテーマなら、テクノロジーの視点から「オンラインでの情報収集は来店のしかたをどう変えたか」のような問いが出てきます。}
\pnote{続けて聞く}{「経済学の2つを、名古屋都市圏に限定して言い換えてください」「12個のうち、大学2年生が半年で調べられそうなものはどれですか。理由も」}
\pnote{確かめる}{候補はあくまで素材です。採点は自分で行い、点数の理由を説明できるものだけを採用します。採用しなかった案と理由をメモしておくと、AI活用記録になります。}

\promptitem{05}{問いを絞る\quad\usedmark}{第2回。候補の問いを3条件で評価するとき}
\begin{promptbox}
次の問いについて、「大学2年生が半年で調べて答えられるか」「範囲は適切か」「答えることに意義があるか」の3点から評価し、必要なら狭める案を示してください。

問い：（候補の問い）
\end{promptbox}
\pnote{解説}{3条件は、授業で示した良い問いの条件そのものです。AIの評価は参考にしますが、点数は自分で付け直します。「意義」の点数は、AIの評価と違ってかまいません。}
\pnote{続けて聞く}{「狭める案を3つ出してください。対象、地域、期間のどれを絞ったかを明示してください」「狭めすぎて意義が薄れていないか、確認してください」}
\pnote{確かめる}{「調べて答えられるか」の最終判断は、実際に文献やデータがあるかどうかで決まります。第3回の検索で確かめます。}

\promptitem{06}{「なぜ気になるのか」から問いを育てる}{第1〜2回。関心のあるテーマと理由しかないとき}
\begin{promptbox}
私は経済学部の2年生です。「（テーマ）」が気になっています。理由は「（なぜ気になるのか。例：実家の近くの自動車ディーラーがいつも空いているから）」です。

この関心を出発点にして、研究の問いの候補を5つ提案してください。各候補には次を添えてください。

1. 問い（疑問文で1行）

2. その問いに答えるために何を調べればよいか（データや資料の種類）

3. 大学2年生が半年で取り組める理由、または難しい理由
\end{promptbox}
\pnote{解説}{第1回の宿題で書いた「なぜ気になるのか」をそのまま使います。理由を渡すと、テーマだけを渡すより、自分の関心に近い問いが出ます。調べ方をセットで出させることで、「調べて答えられるか」を最初から意識できます。}
\pnote{続けて聞く}{「候補2と候補4を組み合わせた問いは作れますか」「候補3の『調べること』は、具体的にどの統計や調査で調べられますか。名前だけ挙げて、数値は挙げないでください」}
\pnote{確かめる}{「調べればよいこと」に挙がった資料の実在は、第3回のやり方で確認します。}

\promptitem{07}{問いの形を整える}{第2回、第5回。「〜について」を疑問文にしたいとき、用語があいまいなとき}
\begin{promptbox}
次の文を、研究の問い（答えを求める疑問文）の形に直してください。

元の文：（例：若者の車離れについて）

条件：

・「なぜ〜なのか」「〜はどのように〜か」「〜は〜にどう影響するか」のいずれかの形にする

・対象（誰）、地域、期間がわかるように具体化する

・あいまいな語（例：「変わった」「影響」）があれば、何を測れば確かめられるかを示す

3通りの案を出し、それぞれの違いを1行で説明してください。
\end{promptbox}
\pnote{解説}{「〜について」は問いではありません。対象、地域、期間を具体化させ、あいまいな語を「何を測れば確かめられるか」に置き換えさせます。第5回で査読者役のAIが指摘する「用語の明確さ」を、前倒しで確認できます。}
\pnote{続けて聞く}{「案2の『来店のしかた』を、回数と目的に分けて書き直してください」}
\pnote{確かめる}{3案のどれを選ぶかは自分で決め、理由を書き残します。}

% ------------------------------------------------------------
\psubsection{先行研究を探す（第3回）}

\promptitem{08}{探す場所の見当をつける\quad\usedmark}{第3回。データベースで検索する前}
\begin{promptbox}
次のリサーチ・クエスチョンに関連する学術分野と、文献検索に使うキーワードを日本語と英語でそれぞれ10個ずつ挙げてください。

文献名や著者名は挙げないでください。

問い：（あなたの問い）
\end{promptbox}
\pnote{解説}{「文献名や著者名は挙げないでください」がこのプロンプトの要です。LLMは文献を捏造しやすいので、頼むのはキーワードまでにします。英語のキーワードがあると、Google Scholar で海外の研究まで探せます。自動車ディーラーへの来店の問いなら、マーケティング（消費者行動論）、流通論、交通経済学、情報行動論などが挙がります。}
\pnote{続けて聞く}{「キーワードを、意味の広いものから狭いものへ並べ替えてください」「隣接する分野をさらに3つ挙げてください」}
\pnote{確かめる}{キーワードは、データベースで実際にヒットしたものだけを使います。検索した日と件数を記録します。}

\promptitem{09}{要旨から関連の深さを判断する\quad\usedmark}{第3回。データベースで見つけた文献を読み分けるとき}
\begin{promptbox}
次の論文の要旨は、私の問い「（あなたの問い）」にどの程度関係しますか。

関係する点と、関係しない点を分けて簡潔に説明してください。

要旨：（データベースで見つけた要旨を貼る）
\end{promptbox}
\pnote{解説}{AIに文献を探させるのではなく、自分が見つけた文献の読み分けを手伝わせる使い方です。関係する点としない点を分けさせることで、その文献を採用した理由、しなかった理由が記録に残ります。}
\pnote{続けて聞く}{「この要旨の方法（対象、期間、データ）は、私の問いの対象とどう違いますか」}
\pnote{確かめる}{要旨を入力するときは、公開範囲と利用条件を確認します。採否の最終判断は、本文を読んでから行います。}

\promptitem{10}{検索式を組み立てる}{第3回。データベースの検索語を決めるとき}
\begin{promptbox}
次の問いについて、CiNii Research と Google Scholar で使う検索式を作ってください。

問い：（あなたの問い）

条件：

・日本語用と英語用を分ける

・同義語や言い換えを OR でつなぎ、異なる概念を AND でつなぐ

・広い式（ヒットが多い）と狭い式（絞り込み）を各2つ

・文献名や著者名は挙げない

各式について、どの概念を表しているかを1行で説明してください。
\end{promptbox}
\pnote{解説}{同義語を OR で、異なる概念を AND でつなぐのが検索式の基本です。「若者」「若年層」「20代」のような言い換えを自分で思いつくのは難しいので、AIに出させます。}
\pnote{続けて聞く}{「ヒットが多すぎました。狭める語を3つ提案してください」「英語の式に、米国と英国で綴りや言い方が違う語はありますか」}
\pnote{確かめる}{検索式の記法はデータベースごとに違います。実際の検索窓の仕様に合わせて手直しし、検索した日と件数を記録します。}

\promptitem{11}{公的統計や白書のあたりをつける}{第3回。数値の出どころを探すとき}
\begin{promptbox}
次の問いに関係しそうな、日本の公的統計、白書、業界団体の調査を挙げてください。

問い：（あなたの問い）

条件：

・実施している省庁や団体の名前と、調査の名前を挙げる

・どのような項目が調べられているかを1行で説明する

・年や数値は挙げない（私が公式サイトで確認します）

・確信がないものには「要確認」と付ける
\end{promptbox}
\pnote{解説}{乗用車市場動向調査、情報通信白書、通信利用動向調査のような調査は、こうして見当をつけてから公式サイトで探します。「年や数値は挙げない」と書くことで、ハルシネーションの数値をレポートに持ち込まずにすみます。}
\pnote{続けて聞く}{「この調査の中で、20代の回答だけを取り出せる項目はありそうですか」}
\pnote{確かめる}{必ず公式サイトで、調査の実在、最新の版、該当ページを確認します。数値は原資料から転記し、版と発行年、URLを文献リストに書きます。}

\promptitem{12}{自分で確認した文献情報を整形する}{第3回、第5回。文献リストを作るとき}
\begin{promptbox}
次の文献情報を、「著者（発行年）題名、掲載誌、巻（号）、頁、DOIまたはURL」の順に整えてください。

条件：

・私が記した情報だけを使い、欠けている項目は「不明」と書く。推測で補わない

・表記ゆれ（全角と半角、スペース、記号）を統一する

文献情報：

1. （データベースで確認した情報を貼る）

2. （同上）
\end{promptbox}
\pnote{解説}{AIに文献を「探させる」のではなく、自分が確認した情報を「整える」のに使います。「推測で補わない」が要です。欠けた項目をAIがもっともらしく埋めてしまうことを防ぎます。}
\pnote{確かめる}{整形後も、年、巻号、頁を原文献と照合します。授業で案内する書式に合わせて最終調整します。}

% ------------------------------------------------------------
\psubsection{読む、批判的に考える（第4回）}

\promptitem{13}{構造をつかむ\quad\usedmark}{第4回。論文を読む前の地図を作るとき}
\begin{promptbox}
次の論文について、「問い」「方法」「主な結果」「著者が認めている限界」をそれぞれ2文以内でまとめてください。書かれていない項目は「記載なし」としてください。

（論文の要旨や該当部分を貼る）
\end{promptbox}
\pnote{解説}{「記載なし」を許すのがポイントです。ないものを作らせません。2文以内と制限することで、要点だけが出ます。この要約は読む前の地図であって、原文の代わりにはなりません。}
\pnote{続けて聞く}{「主な結果について、本文のどの部分（節や表）が根拠になっていますか」「限界について、著者が挙げたものと、あなたが考える限界を分けてください」}
\pnote{確かめる}{原文の頁を書きながら照合し、4つのずれ（問いを広げる、条件を落とす、結果を言い過ぎる、限界を落とす）がないかを確かめます。論文の全文を入力するときは、大学のAI利用規程と資料の利用条件を確認します。}

\promptitem{14}{難しい段落をかみ砕く}{第4回。原文が読みにくいとき}
\begin{promptbox}
次の段落を、経済学部の2年生にわかるように言い換えてください。

条件：

・専門用語はそのまま残し、直後に（　）で短い説明を付ける

・原文にない情報や例は足さない

・原文の主張を弱めたり強めたりしない

・最後に、原文を読むときに注意すべき語句を3つ挙げる

段落：（貼る）
\end{promptbox}
\pnote{解説}{言い換えは、原文を読む手がかりとして使います。「原文にない情報を足さない」「主張を弱めたり強めたりしない」と書くことで、要約で起きやすい言い過ぎを抑えます。専門用語を残させるのは、原文に戻ったときに対応がわかるようにするためです。}
\pnote{続けて聞く}{「『（用語）』は、この分野ではどういう意味で使われますか。一般的な意味との違いも」}
\pnote{確かめる}{言い換えで主張が強まっていないかを原文と比べます。レポートで引用するのは原文からです。}

\promptitem{15}{主張と根拠の対応を確かめる}{第4回。批判的に読むとき}
\begin{promptbox}
次の論文（の一部）について、著者の主張と、その主張を支えている根拠を対応表にしてください。

条件：

・主張ごとに、根拠が「データ（表、図、数値）」「先行研究の引用」「著者の推論」のどれかを示す

・根拠が示されていない主張があれば、「根拠なし」と明示する

・本文に書かれていないことは補わない

（該当部分を貼る）
\end{promptbox}
\pnote{解説}{批判的に読むとは、主張と根拠の対応を確かめることです。著者の推論だけで支えられている主張や、「根拠なし」の主張が、反論の入口になります。}
\pnote{続けて聞く}{「推論だけで支えられている主張のうち、どんなデータがあれば確かめられますか」}
\pnote{確かめる}{表の各行を、原文の頁と表番号で照合します。AIが「根拠なし」とした主張に、実は本文の別の節で根拠が示されていることがあります。}

\promptitem{16}{反論を引き出す\quad\usedmark}{第4回。主張の弱点を探すとき}
\begin{promptbox}
この論文の主張「（主張）」に対して、最も強い反論を3つ挙げてください。

それぞれについて、どのようなデータがあれば検証できるかも示してください。
\end{promptbox}
\pnote{解説}{「どのようなデータがあれば検証できるか」まで出させると、反論が具体的になり、自分の調べ方のヒントにもなります。架空論文Ｘなら、「買った人だけを調べている」「情報を多く集める人はもともと車に詳しいのでは」「情報源の種類を区別していない」のような反論が挙がります。}
\pnote{続けて聞く}{「この反論に、論文の中で答えている箇所はありますか。なければ『なし』と答えてください」。答えが出ても、自分で原文を確かめます。}
\pnote{確かめる}{答えられるかの判定（○△×）は、原文の根拠で行います。答えられない反論は論文の範囲の外であり、自分の問いの出発点になります。}

\promptitem{17}{別の解釈を探す}{第4回。結果の読み方を疑うとき}
\begin{promptbox}
次の論文の結果「（結果。例：情報探索が多い人ほど訪問店舗数が少ない）」について、著者の解釈「（解釈）」以外に成り立ちうる説明を3つ挙げてください。

それぞれについて、次を示してください。

・その説明が正しいとすれば、どんなデータに違いが出るか

・この論文のデータだけで区別できるか
\end{promptbox}
\pnote{解説}{相関と因果の違い、両方に共通する別の要因、対象の偏りなど、同じ結果に別の説明がつくことを確かめる練習です。「車に関心が高い人は探索も多く、店にもよく行く」のような、第三の要因に気づけます。}
\pnote{続けて聞く}{「著者は、この別の説明をどう扱っていますか。統制変数や限界の節に書かれていますか」}
\pnote{確かめる}{著者がすでに統制した変数や、限界の節で認めていることが多いので、原文で確認してから反論にします。}

\promptitem{18}{論文が答えていないことを探す}{第4回。自分の問いとの隙間を見つけるとき}
\begin{promptbox}
次の論文（要旨と結論）が答えていない問いを5つ挙げてください。

条件：

・論文の対象（誰、どこ、いつ）、方法、測った項目の範囲から考える

・それぞれ、論文のどの限定からその問いが残るのかを1行で示す

・私の問い「（あなたの問い）」と関係の深い順に並べる

（要旨と結論を貼る）
\end{promptbox}
\pnote{解説}{第4回のワーク3「未解決の要素」をそのまま頼むプロンプトです。対象、方法、測った項目の「限定」から考えさせると、根拠のある隙間が出てきます。架空論文Ｘなら、来店の「目的」が調べられていないこと、地方と都市の違いが残っていることが、自分の問いにつながります。}
\pnote{続けて聞く}{「1番目の問いに答えるには、論文の方法をどう変える必要がありますか」}
\pnote{確かめる}{「答えていない」と出たものが、実は本文で扱われていることがあります。限界の節と本文で確認します。}

% ------------------------------------------------------------
\psubsection{問いを磨き、確定する（第5回）}

\promptitem{19}{査読者として批判させる\quad\usedmark}{第5回。問いを確定する前}
\begin{promptbox}
あなたは厳しい査読者です。次のリサーチ・クエスチョンについて、

1. 半年で調べて答えられるか

2. 先行研究に対する新しさがあるか

3. 範囲は適切か

4. 用語は明確か

の観点から問題点を指摘し、それぞれ改善案を示してください。

問い：（あなたの問い）

背景：（200字程度）

主な先行研究：（3件）
\end{promptbox}
\pnote{解説}{4つの観点は第5回の基準そのものです。背景と先行研究を渡すことで、指摘が一般論ではなく自分の問いに即したものになります。例えば、「どう変えたのか」という過去との比較が半年では難しいという指摘を受けたら、現状を記述する問いに修正します。}
\pnote{続けて聞く}{「指摘1について、問いを変えずに補強できる方法はありますか」「改善案の問いと私の元の問いを並べて、何が変わったかを表にしてください」}
\pnote{確かめる}{すべての指摘を受け入れる必要はありません。補強か修正かを自分で決め、修正前と修正後の問いを記録します。改善案の文言はそのまま使わず、自分で書き直します。}

\promptitem{20}{反論に答える練習\quad\usedmark}{第5回。指摘に補強で応じるとき}
\begin{promptbox}
私の問いに対して「（AIが挙げた反論）」という批判があります。

この批判に答える論点を3つ挙げ、それぞれどんな根拠が必要かを示してください。

根拠が見つかっていない論点は、そのように明示してください。
\end{promptbox}
\pnote{解説}{「根拠が見つかっていない論点は明示」が要です。根拠があるかのように書かせません。例えば、「新しさがない」という指摘に対しては、先行研究は訪問店舗数を扱い、自分は来店の目的と車の利用が多い地域に注目する、という違いを示して反論できます。}
\pnote{続けて聞く}{「論点2の根拠は、どんな種類の資料（統計、論文、調査）で探せますか。名前だけ挙げて、数値は挙げないでください」}
\pnote{確かめる}{論点の根拠は、手元の文献と公的統計で裏付けます。見つからなければ、問いを修正するか、レポートの限界として書きます。}

\promptitem{21}{調べ方の見通しを立てる}{第5回。レポートの「調べ方の見通し」を書くとき}
\begin{promptbox}
次の問いに答えるために、大学2年生が半年で実行できる調べ方を提案してください。

問い：（確定した問い）

条件：

・「公的統計などの既存データ」「自分で行う調査（アンケート、聞き取り）」「文献」の3つに分けて、それぞれ具体的に

・必要な人数や件数、かかる期間の目安

・その調べ方でわかることと、わからないこと（限界）

・倫理面で注意すべきこと（個人情報、依頼のしかた）
\end{promptbox}
\pnote{解説}{レポートの「調べ方の見通し」の下書きになります。限界を同時に出させておくと、発表で「20人で20代と言えるのか」のような質問を受けたときに、傾向を探る調査と位置づけて答えられます。}
\pnote{続けて聞く}{「アンケートの質問項目の案を5つ。先行研究の分類に合わせてください」「聞き取りを頼むときの依頼文の案を150字で」}
\pnote{確かめる}{人数や期間の目安は根拠のない数字です。授業の期間と自分の時間で判断します。アンケートや聞き取りを行うときは、担当教員に相談します。}

\promptitem{22}{発表の想定問答}{第5回。1人3分の発表の前}
\begin{promptbox}
次の内容で3分の発表をします。聞いている教員や学生が尋ねそうな質問を、厳しい順に5つ挙げてください。各質問に、私が答えるときに押さえるべき論点を1行で添えてください。

問い：（確定した問い）

背景と意義：

先行研究でわかっていること：

調べ方の見通し：
\end{promptbox}
\pnote{解説}{第5回の「1枚にまとめる」4項目をそのまま渡します。想定される質問を事前に知っておくと、答えられないことを「限界」として自分から言えます。}
\pnote{続けて聞く}{「質問1に対する答えを、私の材料だけで60秒以内にまとめるとどうなりますか」}
\pnote{確かめる}{答えの論点はAIの案です。自分の材料で実際に答えられるかを確認し、答えられないものは今後の課題として言う準備をします。}

% ------------------------------------------------------------
\psubsection{レポートに仕上げる}

\promptitem{23}{構成案を一緒に作る}{レポートを書き始めるとき}
\begin{promptbox}
リサーチ・クエスチョン確定レポートの構成案（アウトライン）を作ってください。

必ず含める項目：確定した問いと選んだ理由／背景と意義／先行研究の整理（文献リスト付き）／調べ方の見通し／AI活用記録

私の材料：

・問い：

・背景：

・先行研究（実在を確認したもの）：

・問いの修正の経緯：

条件：

・各節に、書く内容を2〜3行で示す

・私の材料にないことは「材料が必要」と書き、作らない

・文字数の配分の目安を示す
\end{promptbox}
\pnote{解説}{構成案は「一緒に作る」ものです。本文を書くのは自分です。「材料にないことは作らない」と書くことで、AIがもっともらしい背景や文献を埋めてしまうのを防ぎます。}
\pnote{続けて聞く}{「先行研究の節を、文献を年代順に並べる案と、論点別に並べる案の2通りで」}
\pnote{確かめる}{「材料が必要」の箇所を埋めるのは自分です。AIの出力を本文にそのまま貼りません。}

\promptitem{24}{自分の文章を推敲させる}{レポートの下書きができたとき}
\begin{promptbox}
次の文章を推敲してください。ただし書き直しはせず、指摘だけをしてください。

観点：

1. 誤字脱字、表記ゆれ

2. 主語と述語のねじれ、一文が長すぎる箇所

3. 根拠を示さずに断定している箇所

4. 論理の飛躍（前の文からつながらない箇所）

5. 「〜と言われている」「一般に」のような、あいまいな根拠

指摘は、該当箇所を引用し、何が問題かを1行で。

文章：（自分が書いた文章）
\end{promptbox}
\pnote{解説}{「書き直さず指摘だけ」にすることで、文章は自分のものにとどまります。AIに書き直させた文章を提出物にすることは、この授業の「しないこと」にあたります。観点3と5は、出力を使う前のチェックと同じです。}
\pnote{続けて聞く}{「指摘4の箇所について、間に補うべき説明は何ですか。文は書かず、内容だけ示してください」}
\pnote{確かめる}{指摘をどう扱ったか（直した、直さなかった、理由）もAI活用記録に書きます。すべて直す必要はありません。}

\promptitem{25}{AI活用記録の下書きを整える}{対話の終わり。AI活用記録を書くとき}
\begin{promptbox}
この対話の中で私があなたに頼んだことを、AI活用記録の形に整理してください。

形式：表。列は「使用した場面」「主なプロンプト（要約）」「出力をどう扱ったか（採用／修正／不採用）」「検証が必要な事項」

条件：

・「出力をどう扱ったか」は、私がこの対話で明示したものだけを書き、不明なら空欄にする

・「検証が必要な事項」には、この対話で出た固有名詞、数値、文献をすべて挙げる
\end{promptbox}
\pnote{解説}{対話の最後に頼むと、AI活用記録の下書きになります。採用、修正、不採用の判断は自分のものなので、空欄を自分で埋めます。「検証が必要な事項」を挙げさせることで、確認漏れを防げます。}
\pnote{確かめる}{記録は自分の言葉で仕上げます。「検証したこと」の列（何で確かめたか）は自分で追記します。会話が残らない設定のときは、対話ごとに記録します。}

% ------------------------------------------------------------
\psubsection{検証のためのプロンプト（どの場面でも）}

\promptitem{26}{根拠と出典を尋ねる}{出力を使う前}
\begin{promptbox}
今の回答について、次を分けて示してください。

1. 根拠となる資料（調査名、文献、公式サイトなど）。確信がないものには「要確認」と付ける

2. 資料にもとづかない、あなたの推論の部分

3. この回答で私が自分で確認すべき固有名詞と数値の一覧
\end{promptbox}
\pnote{解説}{「根拠は」「出典は」と尋ねるのが検証の第一歩です。資料と推論を分けさせると、どこを確認すべきかがわかります。ただし、挙がった資料そのものが捏造であることもあります。}
\pnote{確かめる}{挙がった資料は必ず自分で開きます。URLが付いていても開いて、内容が書かれているかを見ます。「要確認」と付かなかったものも確認します。}

\promptitem{27}{確かなこと、推測、わからないことに分けさせる}{事実や数値を尋ねるとき}
\begin{promptbox}
「（質問）」について答えてください。ただし、回答を次の3つに分けてください。

A. 確かだと言えること（広く確認されている事実）

B. 推測や一般的な傾向（例外がありうること）

C. わからないこと、あなたの学習データでは判断できないこと（最新の数値、特定の地域の状況など）

Bについては、どんな条件で成り立たなくなるかも示してください。
\end{promptbox}
\pnote{解説}{AIの口調の自信と内容の正しさは無関係です。3つに分けさせると、検証の優先順位がつきます。知識のカットオフがあることを前提に、Cを認めさせるのがねらいです。}
\pnote{確かめる}{Aに分類されたことも検証の対象です。AIが「確か」と言ったことが正しいとは限りません。}

\promptitem{28}{反対の見解と例外を探す}{主張を固める前}
\begin{promptbox}
「（主張）」という主張に反する証拠、反対の見解、例外となる事例を挙げてください。

条件：

・主張を支持する内容は書かない

・それぞれ、どんな資料やデータで確かめられるかを示す

・文献名や著者名は挙げず、「どの分野のどんな研究か」で示す
\end{promptbox}
\pnote{解説}{出力を使う前のチェックの「反対の見解や例外はないか」を、そのまま聞くプロンプトです。AIは利用者の主張に沿った内容を出しやすく、学習データの偏りもあります。反対側を意識的に探させます。}
\pnote{続けて聞く}{「反対の見解のうち、私の問いの対象（名古屋都市圏の20代）にも当てはまるものはどれですか」}
\pnote{確かめる}{反対の見解が「ある」と出ても、資料で確かめてから使います。確かめられたものは、レポートで先行研究の整理や限界として書きます。}

% ============================================================
% 付録B　用語集
% ============================================================
\clearpage
\section{出席確認クイズで扱ったAIの用語}

第1〜5回の出席確認クイズで扱った用語を、クイズの解説にもとづいてまとめました。

\noindent\begin{stab}{colspec={Q[l,wd=36mm,m] X[l,m]},rowsep=3pt}
  用語 & 意味 \\
  AI（人工知能） & Artificial Intelligence。人間の知的な働き（言語理解、推論、学習など）をコンピュータで実現しようとする技術の総称 \\
  生成AI & 学習したパターンをもとに、文章、画像、音声、プログラムなどの新しいコンテンツを作り出すAI。Copilot や ChatGPT、画像生成AIなど \\
  機械学習 & データからパターンや規則性を学び、手順を明示的に書かなくても判断や予測ができるようにする手法。生成AIも機械学習（特に深層学習）にもとづく \\
  教師あり学習、教師なし学習、強化学習 & 正解ラベル付きのデータで入力と正解の対応を学ぶ（教師あり）、正解のないデータからグループ分けや構造を見つける（教師なし）、試行錯誤の結果に報酬を与えて行動を学ばせる（強化学習） \\
  大規模言語モデル（LLM） & Large Language Model。膨大な文章データから言語のパターンを学習したモデル。文脈から次に来る単語（トークン）を確率的に予測して、順に出力することで文章を生成する \\
  トークン & モデルが文章を分割して扱う最小単位。日本語では1〜数文字がひとつのトークンになることが多い。入力と出力の長さもトークン数で数える \\
  コンテキストウィンドウ & モデルが一度の対話で参照できるトークン数の上限。超えると、長い文書の前半が考慮されないなどの問題が起こる \\
  プロンプト & AIに与える指示や質問の文章。質が出力の質を左右する。この授業の4要素は、役割、目的と背景、条件、出力の形式 \\
  temperature & 次の単語を確率分布から選ぶときのランダムさの度合いの設定。同じプロンプトでも出力が毎回少し違う理由 \\
  ハルシネーション & もっともらしいが事実と異なる内容を、自信のある口調で出力すること。存在しない文献や誤った数値など。出力は一次資料で検証する \\
  バイアス & 学習データに含まれる人間社会の偏り（性別、人種、地域など）を、AIが再生産すること。出力を鵜呑みにせず、反対の見解も探す \\
  知識のカットオフ & 学習データがある時点までの情報しか含まず、それ以降の出来事や最新の統計を知らないこと。最新情報は検索連携や一次資料で補う \\
  RAG（検索拡張生成） & 質問に関連する文書や検索結果を取り込み、それを根拠に回答を生成する仕組み。最新情報を扱えるが、読み違いや誤った要約は起こる \\
  AI活用記録 & 提出物に添える記録。使用したツール、使用した場面、主なプロンプト、出力をどう扱ったか、検証したこと、の5項目 \\
\end{stab}

\end{document}
